%% ===========================================================================
%%  Chapter 10 — WASM Emitter
%% ===========================================================================
\section{WASM Emitter}\label{ch:wasm}

\texttt{assembly/sakum\_wasm.s} emits a spec-valid WASM binary
byte-by-byte, proving portable machine-level output.  It compiles
\texttt{func run() -> i32 \{ return 1 + 2 * 3; \}} to a module with magic
\texttt{00 61 73 6d}, version, type\textsection, func\textsection,
export\textsection(\texttt{"run"}), and code\textsection containing
\texttt{i32.const / i32.mul / i32.add / return / end}
(\texttt{41 01 41 02 41 03 6c 6a 0f 0b}).

\begin{lstlisting}[language=SakumAsm,caption={Code-section emit (sakum\_wasm.s:79).}]
    mov dil, 0x0a; call write_byte      # section id 10
    mov dil, 0x0d; call write_byte      # content length 13
    mov dil, 0x01; call write_byte      # func count 1
    mov dil, 0x0b; call write_byte      # body length 11
    mov dil, 0x00; call write_byte      # local decl count 0
    mov dil, 0x41; call write_byte      # i32.const
    mov dil, 0x01; call write_byte
    ...
    mov dil, 0x6c; call write_byte      # i32.mul
    mov dil, 0x6a; call write_byte      # i32.add
    mov dil, 0x0f; call write_byte      # return
    mov dil, 0x0b; call write_byte      # end
\end{lstlisting}

\noindent \texttt{write\_byte} stores one byte at the emit cursor
(\texttt{mov [r15], dil; inc r15}); the buffer is flushed with a raw
\texttt{syscall} write (\texttt{rax = 0x2000004}).

\begin{verbatim}
gcc -arch x86_64 assembly/sakum_wasm.s -o /tmp/wasmgen \
  && /tmp/wasmgen > /tmp/out.wasm && wasm-validate /tmp/out.wasm
\end{verbatim}
